\documentclass[korean,small]{cv}

\setdatewidth{32mm}  % 국문 기간 표기가 길어서 (2022.09 - 2024.06)

% [확인] 아래 표시된 항목은 본인 공개 페이지(GitHub 프로필, til.99jik.com, selab.knu.ac.kr)에서
% 끌어온 추론. 사실 확인 후 주석 지울 것.

\cvname{김정인}
% 자율주행은 뺐음. 읽은 논문에서 추론한 것이고 실적이 없음.
% 실제 실적은 SW 테스팅(KCC 25/26), SIL(ABB), 산업 공정 데이터 ML(UCWIT 25).
\cvtagline{소프트웨어 테스팅 석사과정. LLM을 이용한 테스트 환경 자동 생성과 산업 공정 데이터 기반 예측 모델링}
% 전화번호는 뺐음. 저장소가 public이라 push하면 그대로 노출됨.
% 제출할 때 필요하면 그때 한 줄 넣고 빌드할 것.
% 링크는 둘만. GitHub은 프로젝트 코드 검증 경로라 필수, 개인 사이트는 상위 도메인 하나로 충분.
% til.99jik.com은 99jik.com에서 닿으므로 중복.
% 거주지도 뺐음. 학력과 경력 항목마다 대구가 이미 붙어 있어 중복이고,
% 학술 CV에서 거주지는 관례적으로 안 씀. 회사 이력서 양식에는 따로 적는 칸이 있음.
\cvcontact{%
  \cvmail{99jik@99jik.com}\cvsep
  \cvlink{https://99jik.com}{99jik.com}\cvsep
  \cvlink{https://github.com/99JIK}{github.com/99JIK}%
}

\begin{document}
\makecvheader

% ============================================================ 학력
\section{학력}
% [확인] 지도교수/연구실은 GitHub bio + selab.knu.ac.kr에서 추론
\cventry{2025.03 - 현재}{컴퓨터학부 석사과정}
  {경북대학교}{대구}
  {지도교수: 이우진. 소프트웨어 테스팅 연구실(STLAB). 학위논문 주제 미확정.}

% 성적증명서 기준. 생년월일, 학번, 야간 여부는 CV에 안 넣음.
% 학점은 지역 옆 작은 muted 슬롯으로. 두 학위 모두 표기해 비대칭을 없앰.
% 하나만 적고 하나를 비우면 심사자가 왜 하나만 없는지 알아챔. 숨기는 게 더 불리함.
\cventry{2023.03 - 2024.02}{공학사, 컴퓨터정보공학과 \cvnote{(학사학위 전공심화과정)}}
  {영진전문대학교}{학점 3.50/4.50\cvsep 대구}{}

\cventry{2018.03 - 2023.02}{공업전문학사, 컴퓨터정보계열 웹데이터베이스전공}
  {영진전문대학교}{학점 3.05/4.50\cvsep 대구}{}

% 5년 재학의 이유가 여기서 바로 설명돼야 함. 안 밝히면 심사자가 먼저 묻는다.
% 날짜는 본인 기억 기준. 병적증명서(정부24)로 확인 필요.
\cvline{2018.10 - 2020.06}{대한민국 육군 만기전역}

% ============================================================ 연구 관심 분야
\section{연구 관심 분야}
% 논문 세 편과 ABB 과제에서 실제로 다룬 것만. 자율주행과 퍼징은 실적이 없어 뺐음.
소프트웨어 테스팅, LLM 기반 테스트 생성, SIL 검증 환경, 산업 공정 데이터 기반 예측 모델링

% ============================================================ 논문
\section{논문 \cvsectionnote{\corr 교신저자, 본인 이름 굵게}}

% 최신순으로 위에서 아래로. 번호는 1부터 올라감. 서지사항은 DBpia 등재 기준.
% DBpia 목록에 교신저자가 안 뜨므로 이우진은 따로 붙였음.
\subsection*{국내 학술대회}
\begin{publications}
  \item \me{김정인}, 김덕엽, 서혜령, 이우진\corr.
        대규모 언어 모델을 활용한 임베디드 소프트웨어 시험 환경의 행동 자동 생성.
        \venue{한국정보과학회 학술발표논문집} (KCC 2026), pp. 582--584, 2026.06.
        \cvnote{구현: \cvlink{https://github.com/99JIK/KCC2026}{github.com/99JIK/KCC2026}}

  \item \me{김정인}, 이우진\corr.
        생분해성 섬유 방사 공정의 성능 향상을 위한 샘플링 시점 데이터 활용 연구.
        \venue{유비쿼터스 컴퓨팅 및 웹 정보 기술 학술심포지엄(UCWIT)},
        한국정보과학회 영남지부, 2025.11.

  \item \me{김정인}, 이우진\corr.
        생성형 AI의 테스트케이스 이해 및 활용 능력.
        \venue{한국정보과학회 학술발표논문집} (KCC 2025), pp. 397--399, 2025.07.
\end{publications}

% ============================================================ 특허
% 가출원 상태를 반드시 명시할 것. 정식 출원 전이고 심사도 안 들어갔으므로
% 그냥 "출원"이라 쓰면 과장이 됨. 출원번호 나오면 교체.
% 기여도 수치(30/20/50)는 안 적음. 관례상 발명자 순서만 쓰고, 수치를 적으면 공식 기록상
% 30%로 못박혀 오히려 불리함. 제1발명자라는 사실은 순서로 이미 드러남.
\section{특허}
% 발명자 순서는 출원서 기준으로 본인이 확인해 준 것(김정인 먼저).
% 특허사무소 송부 문서의 대표 발명자란은 이우진이 먼저인데, 그 필드는 연락 책임자를
% 앞세우는 사무소 관행이라 출원서 순서와 다름. 문서 순서로 적었다가 되돌렸음.
% 사무소 관리번호(P-2025-622-02-KR, KL-P-05008)는 출원번호가 아니라 안 씀.
% 정식 출원번호(10-2026-XXXXXXX) 나오면 추가할 것.
\begin{publications}
  \item \me{김정인}, 김덕엽, 이우진. 강건한 최적 공정 변수 조합 선정 기법.
        \cvnote{가출원 완료, 정식 출원 진행 중 (2026.07 기준).
        출원인: 경북대학교 산학협력단. 산업통상자원부/KIAT 과제 P0022335.}
\end{publications}

% ============================================================ 경력
\section{경력}
% 석사 입학은 2025.03이지만 연구실 합류는 2024.12. 사실대로 적으면 공백도 6개월로 줄어듦.
\cventry{2024.12 - 현재}{석사과정 연구원}
  {경북대학교 소프트웨어 테스팅 연구실}{대구}
  {\begin{cvitems}
     \item LLM으로 SIL 테스트 환경의 객체 행위를 자동 생성. 분해-귀납-합성 3단계 프롬프팅으로
           zero-shot 및 few-shot 대비 커버리지 2배 이상 향상. 엘리베이터, 드론, 스마트팩토리
           3개 도메인 환경 객체 11종에서 GPT-4와 Llama-3 70B 모두 일관 (KCC 2026)
     \item 생분해성 섬유 방사 공정 물성 예측에서 샘플링 시점을 시간적 배치 효과의 잠재 변수로
           도입. GBR 기준 $R^2$ 0.724에서 0.811로, RMSE 17.2\% 감소, 산업 허용오차 기준
           정확도 87.1\%에서 92.4\%로 향상. 산업통상자원부/KIAT 과제 P0022335 (UCWIT 2025)
     \item SIL 검증 환경 구축 비용을 낮추는 방법론 연구 (2026, 진행 중)
   \end{cvitems}}

% [확인] 제품명 Ark for CDC는 회사 제품 페이지(iarkdata.com)에서 확인. 담당 제품이 맞는지 확인 필요.
% DBMS 목록은 본인이 직접 분석한 것만. 제품이 지원하는 전체 목록을 옮기면 안 됨.
\cventry{2022.09 - 2024.06}{연구원 \cvnote{(리버스 엔지니어링 담당)}}
  {Arkdata}{대구}
  {\begin{cvitems}
     \item Oracle, PostgreSQL, MariaDB/MySQL, Tibero 등의 이기종 DBMS 트랜잭션 로그
           구조를 리버스 엔지니어링해 CDC(Change Data Capture) 솔루션 개발 지원
     \item Kafka 기반 이기종 DB 간 실시간 데이터 전달 솔루션 설계 및 개발
     \item 3인 연구팀(이사, 과장, 사원)에서 구현 및 분석, 테스트 담당
   \end{cvitems}}

% ============================================================ 강의 경력
% 6개 학기 7건. 프로그래밍 기초를 네 번 맡은 게 재위촉 신호라 굳이 설명 안 해도 읽힘.
\section{강의 경력}
\cventry{2025.03 - 현재}{조교}
  {경북대학교 컴퓨터학부}{대구}
  {\begin{cvitems}
     \item 프로그래밍 기초: 2025-1, 2025 여름, 2025 겨울, 2026 여름 (4개 학기)
     \item 시스템프로그래밍, 창의융합설계: 2025-2
     \item 자바 프로그래밍: 2026-1
   \end{cvitems}}

% ============================================================ 프로젝트
\section{프로젝트}
% ABB는 기업이 아니라 경북대 ICT-ABB 융합전공. A(AI)-B(BigData)-B(Blockchain).
% 지역혁신사업 전자정보 DNA 사업단 소속 교육과정이고, 저장소 README에 "ABB 융합전공" 명시됨.
% 처음에 산학협력 과제로 잘못 적었던 것을 정정함 (ictabb.knu.ac.kr 확인).
% SIL 방법론 연구는 여기서 뺐음. 융합전공이 아니라 연구실 연구라 경력 섹션으로 옮김.
\cventry{2025 - 현재}{ICT-ABB 융합전공 프로젝트}{경북대학교}
  {\cvlink{https://github.com/99JIK/ABB}{github.com/99JIK/ABB}}
  {AU(Action Unit) 기반 실시간 감정 분류 시스템. 최대 30인 동시 처리,
   IoU 기반 다중 얼굴 추적, 감정 변화와 AU 임계값 기반 이벤트 트리거,
   FastAPI REST 및 WebSocket 스트리밍. MediaPipe, TensorFlow.}

% 기간은 저장소 커밋 활동(2026.01 demo - 2026.03)에서 추정. 실제 기간 다르면 수정.
\cventry{2026.01 - 2026.03}{ST영원 스마트 오피스}{외주, 사내 배포 운영 중}
  {\cvlink{https://github.com/99JIK/ST_YOUNGWON}{github.com/99JIK/ST\_YOUNGWON}}
  {\begin{cvitems}
     \item 사내 규정 문서 RAG 질의응답. 규정 문서는 조-항-고정크기 3단 계층 청킹,
           키워드와 벡터를 결합한 하이브리드 검색
     \item OpenAI, Claude, Gemini, 로컬 Ollama를 프로토콜 추상화 뒤에 두어 교체 가능하게 설계.
           유지보수 용이성 요구사항에 대응
     \item FastAPI, ChromaDB, SQLite, Synology NAS(FileStation API) 연동,
           JWT 인증, Docker Compose 배포
   \end{cvitems}}

% 수업 팀 프로젝트지만 내용이 실해서 넣음. ABB 2025년 표정 분석 주제와 결이 같아 서사도 이어짐.
% 기간은 과목코드(2025-2학기)와 저장소 최종 커밋(2025.12)에서 추정.
\cventry{2025.09 - 2025.12}{FOCUS \cvnote{(Facial Observation \& Concentration Understanding System)}}
  {수업 팀 프로젝트}
  {\cvlink{https://github.com/99JIK/FOCUS}{github.com/99JIK/FOCUS}}
  {웹캠 영상으로 학습자의 집중도를 실시간 측정하는 웹 애플리케이션.
   눈 감김, 머리 방향, 시선(두부 회전 보정), 하품, 표정 AU 다섯 지표로 산만도를 점수화.
   MediaPipe Face Landmarker, Vanilla JS, Vite.}

% 원래 COPY_JIKILLER(파이썬 스크립트)였는데 Rust 데스크톱 앱으로 다시 쓰면서 저장소
% 이름도 Provenance로 바뀌었음. 예전 이름은 리다이렉트로 아직 열리지만 같은 이름의
% 저장소가 새로 생기면 깨지므로 현재 이름으로 적음.
% 기간은 저장소 생성(2025.06)과 최종 커밋(2026.08) 기준. 계속 손대고 있음.
\cventry{2025.06--현재}{Provenance \cvnote{(코드 표절 검사)}}
  {개인 프로젝트, Rust}
  {\cvlink{https://github.com/99JIK/Provenance}{github.com/99JIK/Provenance}}
  {텍스트가 아니라 프로그램 구조를 비교해, 변수명을 바꾸거나 주석을 고쳐 써도
   베낀 흔적이 지워지지 않음. Python, Java, C, C++, JavaScript, TypeScript, Go를
   파싱하며 파서가 실행 파일에 내장돼 채점자는 따로 설치할 게 없음.
   제공된 골격 코드는 템플릿으로 지정해 점수에서 제외. 결과는 쌍이 아니라 그룹으로
   묶어 유사도·포함도·밀도로 제시. 크로스플랫폼 데스크톱 앱.
   조교로 맡은 과목의 실기시험 부정행위 적발이 필요해 직접 만듦.}

% TIL 아카이브는 프로젝트에서 뺐음. 자동화 없는 블로그라 프로젝트로는 약하고,
% 링크는 이미 머리말에 있음.

% ============================================================ 기술
\section{보유 기술}
% [확인] GitHub 프로필의 스택 목록 그대로. 실무에서 안 쓴 건 빼는 게 낫다
% 5줄이면 국문이 3쪽으로 넘어가고 마지막 장에 한 줄만 남음. ML을 데이터에 합쳐 4줄로.
% Kafka는 메시징 인프라라 도구 쪽이 더 맞기도 함.
% ML 항목은 UCWIT 2025에서 GBR, LightGBM, RandomForest 실제 사용한 근거 있음.
\cvline{언어}{Python, Java, C/C++, TypeScript, JavaScript}
\cvline{웹/백엔드}{FastAPI, Spring, Node.js, React, Vue.js}
\cvline{데이터/ML}{Oracle, PostgreSQL, MySQL, ChromaDB, scikit-learn, LightGBM}
\cvline{도구}{Docker, Git, Linux, Kafka}

% 수상 섹션은 뺐음. 2012 정보올림피아드 지역예선 동상뿐인데 시점이 14년 전이고 증거가 없어
% 지금 CV에 도움이 안 됨. 나중에 채울 게 생기면 되살릴 것.

\end{document}
